\subsection*{A nonabelian family of refinements and one surviving direction}

\statement{B.59}{Proved}{The residual class-two image lies on one explicit central line}
Put $P=\Gamma_{0,12}$, $\gamma_2P=[P,P]$, and
$\gamma_3P=[\gamma_2P,P]$.  There is an explicitly specified primitive
element $w\in\gamma_2P/\gamma_3P\cong\mathbb Z^{165}$ such that
the image of $\mathcal K$ in the full profinite completion of
$G=P/\gamma_3P$ is contained in
\[
                    \widehat{\mathbb Z}\,w
          \ \subset\ \widehat{\gamma_2P/\gamma_3P}.
                                                               \tag{B.59.1}
\]
The right side is exactly the $A_5$-invariant kernel of the twelve
rigid-seven maps on the completed degree-two quotient.
This statement does not assert that any nonzero point of this line
is realized by $\mathcal K$.

Choose $c\in\gamma_2P$ representing $w$, and define
\[
 R=\langle\gamma_3P,c\rangle\trianglelefteq P,
 \qquad Q=P/R.
                                                               \tag{B.59.2}
\]
Then $Q$ is a finitely generated torsion-free nilpotent group of
class two, with abelianization $\mathbb Z^{54}$ and commutator
subgroup $\mathbb Z^{164}$, and
\[
              \mathcal K\subseteq\overline R.
                                                               \tag{B.59.3}
\]
For all finite-index subgroups $L_1,L_2\leq P$ containing $R$,
the product $(H_T\cap L_1)(H_5\cap L_2)$ is closed in the full
profinite topology of $P$.  Thus every finite quotient factoring
through this particular nonabelian class-two quotient $Q$
gives valid refinements of B.52.

\begin{proof}
\emph{An integral presentation in degree two.}
Let $I=\{1,\ldots,12\}\setminus\{2\}$, in increasing order,
and use pure Artin generators $A_{ij}$ for $i,j\in I$.
The identification of $P$ with $PB_{11}$ modulo its central
full twist removes the single degree-one relation
$\sum_{i<j}A_{ij}=0$ and does not change degree two.
The integral infinitesimal braid relations give the basis
\[
 t_{ijk}=[A_{ij},A_{ik}],\qquad i<j<k\text{ in }I,
                                                               \tag{B.59.4}
\]
of $\gamma_2P/\gamma_3P$; there are $\binom{11}{3}=165$
such triples.  Here the group commutator is $aba^{-1}b^{-1}$.
The other brackets on a triple are
$[A_{ij},A_{jk}]=-t_{ijk}$ and
$[A_{ik},A_{jk}]=t_{ijk}$; brackets on disjoint pairs vanish.
These integral presentations and the central direct-product
splitting are recorded in
\citet[Theorem 2.2(ii), equations (3.1)--(3.2), and
Theorem 3.1(ii)]{DE_KaabiVershinin2012}.
Independence in degree two can also be checked by forgetting
all but each triple of disk strands: the resulting $PB_3$
degree-two quotient is infinite cyclic and detects its own
$t_{ijk}$ alone.

For vectors $v,u$ in the degree-one coordinates of B.58,
the coefficient of $t_{ijk}$ in their bracket is
\[
 v_{ij}u_{ik}-v_{ik}u_{ij}
 -v_{ij}u_{jk}+v_{jk}u_{ij}
 +v_{ik}u_{jk}-v_{jk}u_{ik}.                              \tag{B.59.5}
\]
Represent these vectors using zero coefficient on the omitted
constant pair.  For the source omitted generator $A_{46}$,
its degree-one vector is minus the sum of the other $54$
generators.  The full-twist vector is central, so the bracket
is independent of the chosen representative.
Apply (B.59.5) to the integral degree-one matrix from B.58.
Each seven-point target contributes $\binom63=20$ rows.
This constructs an exact $240\times165$ integer matrix $C$.

The program \nolinkurl{check_ambient_rigid7_class2.py}
checks the convention on the identity map, constructs $C$,
and computes
\[
 \operatorname{rank}C=159,\qquad
 \operatorname{Smith}(C)=(1^{159},0^6).
                                                               \tag{B.59.6}
\]
\begin{sloppypar}
A supplied $159\times159$ minor has determinant $-1$;
its integer inverse is verified in both multiplication orders.
The complete matrix and its ordered row and column labels are in
\nolinkurl{ambient_rigid7_class2_certificate.json}.
The six nullspace basis columns supplied there form an
integer matrix $N$ with $CN=0$.
The second certificate selects a $6\times6$ minor of $N$
of determinant $\pm1$ and supplies its integer inverse.
Consequently $N\mathbb Z^6$ is the full integral kernel.
In particular, ambient degree-two detection is not injective.
\end{sloppypar}

\emph{The actual finite-group action.}
Use the order-two and order-three generators of $T$, followed
by the chosen order-five generator $D$, recorded as labelled
permutations in \nolinkurl{rigid7_subsets.json}.
They generate $T$ of order $12$ and $A_5$ of order $60$.

For a permutation $\sigma$, the degree-one
action is obtained by taking, for each pair $(i,j)$ of the
ambient basis, the column of B.58's matrix $A$ indexed by
$\{\sigma(i),\sigma(j)\}$.
Equation (B.59.5) then gives its degree-two action.
Inner automorphisms by pure mapping classes are trivial
on both these graded quotients, so this calculation is the
action of the actual geometric $A_5$.

The script \nolinkurl{check_rigid7_class2_invariants.py}
verifies that these three actions preserve $N\mathbb Z^6$,
and gives the following matrices in that ordered basis:
\[
 Q_2=\begin{pmatrix}
 0&0&1&0&0&-1\\0&0&0&1&0&0\\1&0&0&0&0&1\\
 0&1&0&0&0&0\\0&0&0&0&1&0\\0&0&0&0&0&1
 \end{pmatrix},\qquad
 Q_3=\begin{pmatrix}
 0&0&0&0&-1&0\\1&0&0&0&0&0\\0&0&0&1&0&1\\
 0&0&1&0&0&-1\\0&-1&0&0&0&0\\0&0&0&1&0&0
 \end{pmatrix},
\]
\[
 Q_5=\begin{pmatrix}
 1&0&0&0&0&0\\0&0&0&0&-1&0\\0&1&0&1&0&0\\
 0&0&1&0&0&-1\\0&0&0&0&0&-1\\0&0&0&1&0&0
 \end{pmatrix}.
                                                               \tag{B.59.7}
\]
The stacked matrix of $Q_2-1,Q_3-1,Q_5-1$ has Smith
factors $(1,1,1,1,1,0)$.  Its kernel is generated by
\[
          v=(1,1,2,1,-1,1)^{\mathsf T},\qquad w=Nv.
                                                               \tag{B.59.8}
\]
This specifies $w$ completely; its $165$ integer coefficients
in the basis (B.59.4), including its $45$ nonzero coefficients,
are also supplied in
\nolinkurl{rigid7_class2_invariants_certificate.json}.
Some coefficient equals $1$, so $w$ is primitive.
One explicit choice in (B.59.2) is therefore the lexicographically
ordered product $c=\prod_{i<j<k}[A_{ij},A_{ik}]^{w_{ijk}}$.
Both programs' actual standard outputs are retained in their
corresponding \nolinkurl{.out} files.
The independent root certificate
\nolinkurl{code/verify_B59_certificate.py}
checks $CN=0$, the primitive kernel minor, ranks $159$ and
$5$ modulo $101$, primitivity of $w$, and the consistency
of the resulting $60$-element matrix action; its actual
standard output is retained.

These integer kernel calculations remain valid over
$\widehat{\mathbb Z}$: this is a torsion-free, hence flat,
$\mathbb Z$-module.  Equivalently, use the supplied primitive
kernel bases and the nonzero integer minors to solve the
equations; multiplication by every nonzero integer is
injective on $\widehat{\mathbb Z}$.
No averaging by $60$ occurs, and the conclusions include
the primes $2,3,5$.

\emph{Passage to full profinite quotients.}
There is no implicit replacement of a full topology by a
nilpotent topology here.  The group $G$ has an exact central
extension with quotient $\mathbb Z^{54}$ and kernel
$\mathbb Z^{165}$.  Choosing ordered lifts of the first
basis and the central basis gives unique coordinates
$\mathbb Z^{54}\times\mathbb Z^{165}$.
Collection of commutators makes their multiplication an
integer bilinear formula.  Reducing all coordinates modulo
$m$ consequently defines a finite group quotient.
These quotients are cofinal among all finite quotients:
if a finite quotient has exponent $d$, every ordered
normal form with all coordinates divisible by $d$ maps
to its identity.  Thus the full completion has the
corresponding $\widehat{\mathbb Z}$ coordinates, and its
kernel over completed abelianization is precisely
$\widehat{\mathbb Z}^{165}$.
The same argument applies to every target class-two
quotient and to $Q$, whose central lattice is the quotient
by the primitive vector $w$.

By B.58, the image of $\mathcal K$ in $\widehat G$ lies
in that completed central lattice.  B.32 kills its images
in all twelve rigid targets.  Every element of
$\mathcal K=C_{\widehat P}(A_5)$ is fixed by $A_5$.
Equations (B.59.6)--(B.59.8) therefore prove (B.59.1).
Quotienting by $\widehat{\mathbb Z}w$ proves that its image
in $\widehat Q$ is trivial.  The universal property of
full profinite completion identifies
$\ker(\widehat P\to\widehat{P/R})$ with $\overline R$,
which proves (B.59.3).
The factorization argument in B.58 now applies verbatim
with $R$ in place of $[P,P]$, giving the stated products.

Finally $Q$ has a nonzero free abelian commutator subgroup
of rank $164$.  A nonzero commutator has nonzero image in
some finite coordinate quotient, proving that this family
includes genuinely nonabelian finite quotients.
\end{proof}

The single central line in (B.59.1) is a remaining possible
image, not a constructed element of $\mathcal K$.
In particular the argument proves neither
$\mathcal K\subseteq\overline{\gamma_3P}$ nor closure of
all finite-index refinements, even among all class-two
quotients before this one line is removed.
